% ======================================================================
% sec-abstract.tex -- abstract (lead instantiation chosen by \ifleaddatalog)
% ======================================================================
\begin{abstract}
  Many data systems admit multiple admissible outcomes for the same
  input: a logic program may have several stable models; concurrent
  transactions may serialize in one of many orders.
  Classical data provenance cannot even pose its question in such
  settings---it explains how a result was derived, but only after
  something has chosen \emph{which} result to produce.
  We introduce \emph{determination provenance} to track the
  commitments that resolve this ambiguity.
  Modeling resolution as a sequence of irrevocable, state-indexed
  \emph{commitments} drawn from a basis, each observable acquires a
  \emph{support}: the set of complete resolutions---\emph{determinations}---under
  which it holds.
  Supports form the Boolean shadow of a product provenance semiring, and
  the layering of commitments induces a \emph{filtration} measuring each
  observable's \emph{support depth}: how many layers of semantic
  resolution it depends on.
  The result is an algebra of provenance queries that reaches well beyond
  why/why-not: how robustly a tuple holds, at what depth its fate is sealed,
  how its explanation shifts across specifications or semantics, and which
  commitment is most responsible. Discharging a non-monotone step into a
  commitment also re-enables ordinary provenance semirings on negation, so
  counting, probability, security, and cost queries apply to negation-licensed
  facts.
  The filtration is respected by provenance composition, enabling
  compositional robustness analysis and quantitative diagnosis of
  resolution cost.
  \ifleaddatalog
  Our lead instantiation is $\mbox{Datalog}^\neg$: for programs with
  stable models, the classical negation semantics (stratified,
  well-founded, stable) are different reading depths of a single shared
  filtration, and sealing commitments reproduce the support behavior of
  semirings with monus while also handling unstratified negation.
  A second instantiation (transactional isolation) appears in an
  appendix.
  \else
  Our lead instantiation is transactional isolation: isolation levels
  (read committed, snapshot isolation, serializability) are different
  resolving subsets of one shared filtration, with determination depth
  $\Theta(n)$ under scheduling discretion.
  A second instantiation ($\mbox{Datalog}^\neg$ negation semantics)
  appears in an appendix.
  \fi
\end{abstract}
